\section*{Capítulo \ref{ch:b2:genome-diversification} --- Mutaciones y diversificación de los genomas}

\begin{solution}{exo:b2:genome-diversification:1}
GAG: silenciosa (sigue siendo glutamato), una transición (A$\to$G). GTA:
de sentido erróneo (valina), una transversión (A$\to$T). TAA: sin
sentido (terminación), una transversión (G$\to$T). Insertar una base:
un desplazamiento del marco, que reescribe todos los codones siguientes.
\end{solution}

\begin{solution}{exo:b2:genome-diversification:2}
$U = 5\times 10^{-10}\times 4.6\times 10^{6} = 2.3\times 10^{-3}$: una
división de cada 430 produce una \omterm{def:b2:genome-diversification:mutation}{mutación}. En $10^{9}$ células:
$2.3\times 10^{6}$ \omterm{def:b2:genome-diversification:mutation}{mutaciones} nuevas por generación.
\end{solution}

\begin{solution}{exo:b2:genome-diversification:3}
\omterm{def:b2:genome-diversification:hgt}{Transformación}: captación de ADN libre procedente de células lisadas;
requiere un receptor competente; transfiere fragmentos cromosómicos
(genes de la cápsula en los neumococos). Conjugación: contacto entre
células a través de un pilus; requiere un \omterm{def:b2:unicellular-diversity:prokaryote}{plásmido} conjugativo en el
donante; transfiere \omterm{def:b2:unicellular-diversity:prokaryote}{plásmidos} (genes de resistencia) y, desde una cepa
Hfr, genes cromosómicos. \omterm{def:b2:genome-diversification:hgt}{Transducción}: un bacteriófago que ha empaquetado
ADN del hospedador; transfiere un fragmento al azar (generalizada) o los
genes vecinos del sitio del profago (especializada).
\end{solution}

\begin{solution}{exo:b2:genome-diversification:4}
\omterm{def:b2:genome-diversification:polyploidy}{Autopoliploide}: juegos adicionales de una misma especie, a partir de un
gameto diploide (la patata tetraploide, muchos tréboles
autotetraploides). \omterm{def:b2:genome-diversification:polyploidy}{Alopoliploide}: los juegos de dos especies combinados
por hibridación y después duplicados (trigo harinero, algodón, tabaco).
El híbrido diploide AB no tiene dos cromosomas iguales, de modo que
ninguno puede aparearse en la meiosis I; los cromosomas se reparten al
azar y los gametos quedan desequilibrados e inviables.
\end{solution}

\begin{solution}{exo:b2:genome-diversification:5}
$2\times 3.2\times 10^{9}\times 1.2\times 10^{-8} = 77$ \omterm{def:b2:genome-diversification:mutation}{mutaciones} nuevas.
En secuencia codificante: $77\times 0.015 = 1.2$; que cambien un
aminoácido: $1.2\times 0.7 = 0.8$ --- la mayoría de los hijos lleva
alrededor de una sustitución de aminoácido nueva.
\end{solution}

\begin{solution}{exo:b2:genome-diversification:6}
Sucesos: $m(N - N_0) = 2\times 10^{-8}\times 10^{9} = 20$. Células
mutantes: $mN\ln(N/N_0) = 20\times \ln 10^{6} = 20\times 13.8 = 276$.
$p_0 = e^{-20} \approx 2\times 10^{-9}$: todos los cultivos tienen
mutantes, así que contar los cultivos sin ninguno no mide nada. Hace
falta $m(N -
N_0) \approx 1$, es decir, $N \approx 5\times 10^{7}$, donde $p_0 \approx
0.37$ y unas pocas decenas de cultivos dan $m$ dentro de un factor dos.
\end{solution}

\begin{solution}{exo:b2:genome-diversification:7}
El uracilo no tiene cabida en el ADN, de modo que una uracil-ADN
glucosilasa puede retirar sin ambigüedad todo uracilo que encuentre. La
timina es una base normal: tras la desaminación de la 5-metilcitosina se
detecta el desapareamiento G$\cdot$T, pero la maquinaria de reparación no
tiene manera de saber qué hebra es la errónea y la mitad de las veces
``corrige'' la G, fijando una transición C$\to$T. Por eso los sitios CpG
metilados mutan unas diez veces más deprisa que los demás, y a lo largo
del tiempo evolutivo se han convertido en TpG y CpA: los genomas de los
mamíferos solo contienen alrededor de una quinta parte de los CpG que
predice su composición de bases.
\end{solution}

\begin{solution}{exo:b2:genome-diversification:8}
\qty{46}{kb} por minuto. Posiciones: 368, 782, 1150 y
\qty{1840}{kb}; distancias de 414, 368 y \qty{690}{kb}.
\end{solution}

\begin{solution}{exo:b2:genome-diversification:9}
La segunda copia de un cromosoma se aparea con la primera copia del
otro; un entrecruzamiento entre ellas da una cromátida que lleva la
copia~1, la copia~2 y la copia~2$'$ (tres genes) y otra que solo lleva la
copia~1$'$. Una persona con un cromosoma de un solo gen de cada
progenitor tiene dos genes $\alpha$ en lugar de cuatro ($-\alpha/-\alpha$):
las cadenas $\alpha$ se fabrican a la mitad de ritmo, los glóbulos rojos
son pequeños y pálidos, con una anemia leve o nula --- el rasgo
$\alpha$-talasémico.
\end{solution}

\begin{solution}{exo:b2:genome-diversification:10}
Gametos: $n = 21$ (un juego de A, uno de B y uno de D). Híbridos: AB, 14
cromosomas; ABD, 21. Trigo harinero $\times$ almidonero: AABBD, 35
cromosomas --- los juegos A y B se aparean, el juego D no tiene pareja, y
la planta es en gran medida estéril. Cada duplicación dio a cada
cromosoma un homólogo idéntico, de modo que se forman bivalentes y los
gametos quedan equilibrados; el retrocruzamiento con un progenitor produce
plantas con un juego sin pareja, cuyos gametos están desequilibrados, de
modo que el hexaploide está aislado reproductivamente de sus antepasados
--- una especie en un solo paso.
\end{solution}

\begin{solution}{exo:b2:genome-diversification:11}
Crecimiento logístico: $R(t) = N/\bigl(1 + (N - 1)e^{-\gamma Nt}\bigr)$.
La mitad de la población cuando $(N - 1)e^{-\gamma Nt} = 1$: $t = \ln(N -
1)/\gamma N = 20.7/0.5 = \qty{41}{h}$. Sin el antibiótico, las células
sin \omterm{def:b2:unicellular-diversity:prokaryote}{plásmido} crecen un \qty{5}{\%} más deprisa, de modo que la razón entre
portadoras y no portadoras decae como $e^{-st}$ con $s = 0.05\ln 2 =
\qty{0.035}{h^{-1}}$ por hora (\omterm{def:b2:unicellular-diversity:division}{tiempo de generación} de \qty{1}{h}): pasar
del \qty{99}{\%} al \qty{1}{\%} lleva $\ln(99^2)/0.035 = \qty{260}{h}$,
once días --- y la transferencia continua puede mantener el \omterm{def:b2:unicellular-diversity:prokaryote}{plásmido}
indefinidamente a baja frecuencia, listo para el próximo tratamiento con
el fármaco.
\end{solution}

\begin{solution}{exo:b2:genome-diversification:12}
El \omterm{thm:b2:genome-diversification:luria}{test de fluctuación} mostró que las \omterm{def:b2:genome-diversification:mutation}{mutaciones} de resistencia surgen
antes y en ausencia del agente selectivo, en momentos aleatorios, de modo
que la \omterm{def:b2:genome-diversification:mutation}{mutación} es ciega a lo que resultaría útil; la siembra en réplica
mostró lo mismo sin exponer en absoluto las células seleccionadas. La
evolución no es aleatoria porque la selección criba esas variantes ciegas
según sus consecuencias, de forma sistemática y acumulativa. El único
punto en que el propio azar está moldeado es la tasa: los linajes con
demasiados o demasiado pocos errores pierden, de modo que la \omterm{thm:b2:genome-diversification:rate}{tasa de
mutación} es un compromiso adaptado --- pero lo que hace cada \omterm{def:b2:genome-diversification:mutation}{mutación}
sigue siendo impredecible y no elegido.
\end{solution}

\begin{solution}{pb:b2:genome-diversification:1}
\textbf{1.} Suma 119, media $5.95$; media de los cuadrados $11489/20 = 574$;
varianza $574 - 35 = 540$.
\textbf{2.} Suma 324, media $16.2$; media de los cuadrados $5462/20 = 273$;
varianza $273 - 262 = 11$.
\textbf{3.} Las muestras de un único cultivo son de Poisson (varianza
próxima a la media): son extracciones aleatorias de una misma fracción
fija de mutantes. Los cultivos independientes tienen una varianza noventa
veces superior a la media: sus mutantes surgieron en momentos distintos y
aleatorios durante el crecimiento, y cuanto más temprana la \omterm{def:b2:genome-diversification:mutation}{mutación},
mayor su clon.
\textbf{4.} Un premio gordo: una \omterm{def:b2:genome-diversification:mutation}{mutación} en una de las primeras
divisiones, cuyo clon creció después con el cultivo hasta un centenar de
células.
\textbf{5.} Catorce de veinte: $p_0 = 0.70$; $m(N - N_0) = -\ln 0.70
= 0.36$; $m = 0.36/2\times 10^{8} = \qty{1.8e-9}{}$ por división.
\textbf{6.} $1.8\times 10^{-9}\times 2\times 10^{8}\times \ln(2\times
10^{6}) = 0.36\times 14.5 = 5.2$ células mutantes; observado, $5.95$ ---
concordancia dentro del ruido de veinte cultivos.
\textbf{7.} Ninguna muestra salió vacía (todas comparten los mutantes del
cultivo original), de modo que $p_0 = 0$ no da ninguna estimación; las
muestras miden la fracción mutante de un solo cultivo, que depende de
cuándo ocurrieron sus \omterm{def:b2:genome-diversification:mutation}{mutaciones}.
\textbf{8.} $u = m/20 = \qty{9e-11}{}$ por par de bases y por replicación.
\textbf{9.} $U = 9\times 10^{-11}\times 4.6\times 10^{6} =
\qty{4.1e-4}{}$.
\textbf{10.} $10^{12}$ células $\times$ $4.1\times 10^{-4}$ $= 4.1\times
10^{8}$ \omterm{def:b2:genome-diversification:mutation}{mutaciones} nuevas.
\textbf{11.} $3\times 4.6\times 10^{6} = 1.4\times 10^{7}$ cambios
posibles; cada uno surge $4.1\times 10^{8}/1.4\times 10^{7} \approx 30$
veces.
\textbf{12.} Todo cambio de una sola base, incluido todo aquel que
confiere resistencia a un fármaco vencido por una sola \omterm{def:b2:genome-diversification:mutation}{mutación}, está
presente en decenas de células antes de añadir el fármaco; el fármaco
se limita a seleccionarlos.
\textbf{13.} Dos \omterm{def:b2:genome-diversification:mutation}{mutaciones} independientes en la misma célula ocurren con
$m^2 \approx (2\times 10^{-9})^2 = 4\times 10^{-18}$ por división: en
$10^{12}$ células, $4\times 10^{-6}$ por generación --- una célula
doblemente resistente prácticamente nunca preexiste, y por eso la
tuberculosis se trata con varios fármacos a la vez.
\textbf{14.} Con $a = N - 1$ y $E = e^{-\gamma Nt}$: $R = N/(1 +
aE)$, $\mathrm{d}R/\mathrm{d}t = \gamma N\cdot NaE/(1 + aE)^2$; y
$\gamma R(N - R) = \gamma\,\frac{N}{1 + aE}\cdot\frac{NaE}{1 + aE}$,
lo mismo; $R(0) = N/(1 + N - 1) = 1$.
\textbf{15.} $(N - 1)E = 1$: $t = \ln(N - 1)/\gamma N = 20.7/0.5 =
\qty{41}{h}$.
\textbf{16.} $1 + aE = 1/0.99$, $aE = 0.0101$: $t = \ln(10^{9}/0.0101)
/0.5 = (20.7 + 4.6)/0.5 = \qty{51}{h}$.
\textbf{17.} $\gamma (N/2)^2 = \gamma N\cdot N/4 = 0.5\times 2.5\times
10^{8} = 1.25\times 10^{8}$ transferencias por hora.
\textbf{18.} Tasas de crecimiento de $\ln 2 = \qty{0.693}{h^{-1}}$ y
\qty{0.762}{h^{-1}}; la razón mutante/sensible crece como $e^{st}$
con $s = \qty{0.069}{h^{-1}}$; pasar de $1/N$ a 1 lleva $\ln N/s =
20.7/0.069 = \qty{300}{h}$, doce días.
\textbf{19.} La transferencia es proporcional al producto de portadoras y
no portadoras y avanza al ritmo de los encuentros, no de las divisiones;
la descendencia está limitada por la ventaja de crecimiento del mutante,
que es pequeña. Un \omterm{def:b2:unicellular-diversity:prokaryote}{plásmido} atraviesa la población en dos días; un
mutante favorecido, en dos semanas.
\textbf{20.} $U = 4.1\times 10^{-2}$; perjudiciales: $0.041\times 0.88\times
0.3 = 0.011$ por generación.
\textbf{21.} Media de $2.2$: $e^{-2.2} = 0.11$, un descendiente de cada
nueve indemne.
\textbf{22.} Tipo silvestre: $1.1\times 10^{-4}$ por generación, $0.022$ en
200: $e^{-0.022} = 0.98$.
\textbf{23.} Media de $520$ células mutantes; $p_0 = e^{-36} \approx
10^{-16}$: ningún cultivo vacío.
\textbf{24.} Bajo una selección intensa y cambiante (la rotación de
antibióticos de un hospital), la oferta centuplicada de variantes nuevas
compensa la carga, y los mutadores viajan a remolque de las resistencias
que generan. En un ambiente estable no hay nada que ganar y sí una
\omterm{def:b2:genome-diversification:mutation}{mutación} perjudicial cada cien generaciones que perder, de modo que los
linajes mutadores se eliminan.
\textbf{25.} $m = \qty{1.8e-9}{}$ por división hacia la resistencia, $u =
\qty{9e-11}{}$ por par de bases y por replicación, $U = \qty{4.1e-4}{}$
\omterm{def:b2:genome-diversification:mutation}{mutaciones} por genoma y por replicación; el \omterm{def:b2:unicellular-diversity:prokaryote}{plásmido} alcanza la mitad del
intestino en \qty{41}{h} y el \qty{99}{\%} en \qty{51}{h}.
\end{solution}
